\documentclass{amsart}
% For dealing with references we use the comment environment
\usepackage{verbatim}
\newenvironment{reference}{\comment}{\endcomment}
%\newenvironment{reference}{}{}
\newenvironment{slogan}{\comment}{\endcomment}
\newenvironment{history}{\comment}{\endcomment}

% For commutative diagrams we use Xy-pic
\usepackage[all]{xy}

% We use 2cell for 2-commutative diagrams.
\xyoption{2cell}
\UseAllTwocells

% We use multicol for the list of chapters between chapters
\usepackage{multicol}

% This is generally recommended for better output
\usepackage{lmodern}
\usepackage[T1]{fontenc}

% For cross-file-references

% Package for hypertext links:
\usepackage{hyperref}

% For any local file, say "hello.tex" you want to link to please

% Theorem environments.
%
\theoremstyle{plain}
\newtheorem{theorem}[subsection]{Theorem}
\newtheorem{proposition}[subsection]{Proposition}
\newtheorem{lemma}[subsection]{Lemma}

\theoremstyle{definition}
\newtheorem{definition}[subsection]{Definition}
\newtheorem{example}[subsection]{Example}
\newtheorem{exercise}[subsection]{Exercise}
\newtheorem{situation}[subsection]{Situation}

\theoremstyle{remark}
\newtheorem{remark}[subsection]{Remark}
\newtheorem{remarks}[subsection]{Remarks}

\numberwithin{equation}{subsection}

% Macros
%
\def\lim{\mathop{\mathrm{lim}}\nolimits}
\def\colim{\mathop{\mathrm{colim}}\nolimits}
\def\Spec{\mathop{\mathrm{Spec}}}
\def\Hom{\mathop{\mathrm{Hom}}\nolimits}
\def\Ext{\mathop{\mathrm{Ext}}\nolimits}
\def\SheafHom{\mathop{\mathcal{H}\!\mathit{om}}\nolimits}
\def\SheafExt{\mathop{\mathcal{E}\!\mathit{xt}}\nolimits}
\def\Sch{\mathit{Sch}}
\def\Mor{\mathop{\mathrm{Mor}}\nolimits}
\def\Ob{\mathop{\mathrm{Ob}}\nolimits}
\def\Sh{\mathop{\mathit{Sh}}\nolimits}
\def\NL{\mathop{N\!L}\nolimits}
\def\CH{\mathop{\mathrm{CH}}\nolimits}
\def\proetale{{pro\text{-}\acute{e}tale}}
\def\etale{{\acute{e}tale}}
\def\QCoh{\mathit{QCoh}}
\def\Ker{\mathop{\mathrm{Ker}}}
\def\Im{\mathop{\mathrm{Im}}}
\def\Coker{\mathop{\mathrm{Coker}}}
\def\Coim{\mathop{\mathrm{Coim}}}

% Boxtimes
%
\DeclareMathSymbol{\boxtimes}{\mathbin}{AMSa}{"02}

%
% Macros for moduli stacks/spaces
%
\def\QCohstack{\mathcal{QC}\!\mathit{oh}}
\def\Cohstack{\mathcal{C}\!\mathit{oh}}
\def\Spacesstack{\mathcal{S}\!\mathit{paces}}
\def\Quotfunctor{\mathrm{Quot}}
\def\Hilbfunctor{\mathrm{Hilb}}
\def\Curvesstack{\mathcal{C}\!\mathit{urves}}
\def\Polarizedstack{\mathcal{P}\!\mathit{olarized}}
\def\Complexesstack{\mathcal{C}\!\mathit{omplexes}}
% \Pic is the operator that assigns to X its picard group, usage \Pic(X)
% \Picardstack_{X/B} denotes the Picard stack of X over B
% \Picardfunctor_{X/B} denotes the Picard functor of X over B
\def\Pic{\mathop{\mathrm{Pic}}\nolimits}
\def\Picardstack{\mathcal{P}\!\mathit{ic}}
\def\Picardfunctor{\mathrm{Pic}}
\def\Deformationcategory{\mathcal{D}\!\mathit{ef}}

\setlength{\emergencystretch}{3em}
\begin{document}
\title[Stacks Project, Tag 07JG]{332 --- Crystals on affine crystalline sites yield complete modules with topologically quasi-nilpotent integrable connections}
\maketitle
\noindent Copyright (C) 2005--2025 Johan de Jong. Stacks Project authors. Source: \href{https://stacks.math.columbia.edu/tag/07JG}{Tag 07JG}.

\noindent Editorial difficulty note (not author text). Difficulty H2: Beyond extracting a connection, the proof must reconstruct a complete module from its reductions and prove topological quasi-nilpotence. It expands the crystal transport in divided-power diagonal coordinates and uses the triple-diagonal cocycle to identify every coefficient with iterated commuting connection operators, forcing p-adic decay..

\noindent Editorial provenance note (not author text). History: excerpt from revision \texttt{a0444\allowbreak 6e57e\allowbreak c1fbc\allowbreak 252a8\allowbreak 71afc\allowbreak ec775\allowbreak 2fb28\allowbreak 07b14}, crystalline.tex, lines 2996--3106. Reference presentation was adapted; original-author context and core support blocks were retained or added. The mathematical wording is unchanged. Licensed under GNU FDL 1.2 or later; no invariant sections or cover texts. See ../licenses/COPYING. Context blocks reproduce author definitions and supporting results; they are not additional counted problems.

\section{Crystals in quasi-coherent modules}
\label{section-quasi-coherent-crystals}

\noindent
In Situation \ref{situation-affine}.
Set $X = \Spec(C)$ and $S = \Spec(A)$. We are going to
classify crystals in quasi-coherent modules on $\text{Cris}(X/S)$.
Before we do so we fix some notation.

\medskip\noindent
Choose a polynomial ring $P = A[x_i]$ over $A$ and a surjection $P \to C$
of $A$-algebras with kernel $J = \Ker(P \to C)$. Set
\begin{equation}
\label{equation-D}
D = \lim_e D_{P, \gamma}(J) / p^eD_{P, \gamma}(J)
\end{equation}
for the $p$-adically completed divided power envelope.
This ring comes with a divided power ideal $\bar J$ and divided power
structure $\bar \gamma$, see Lemma \href{https://stacks.math.columbia.edu/tag/07KG}{lemma list properties (Tag 07KG)}.
Set $D_e = D/p^eD$ and denote $\bar J_e$ the image of $\bar J$ in $D_e$.
We will use the short hand
\begin{equation}
\label{equation-omega-D}
\Omega_D = \lim_e \Omega_{D_e/A, \bar\gamma} =
\lim_e \Omega_{D/A, \bar\gamma}/p^e\Omega_{D/A, \bar\gamma}
\end{equation}
for the $p$-adic completion of the module of divided power differentials,
see Lemma \href{https://stacks.math.columbia.edu/tag/07KK}{lemma differentials completion (Tag 07KK)}.
It is also the $p$-adic completion of
$\Omega_{D_{P, \gamma}(J)/A, \bar\gamma}$
which is free on $\text{d}x_i$, see
Lemma \href{https://stacks.math.columbia.edu/tag/07HW}{lemma module differentials divided power envelope (Tag 07HW)}.
Hence any element of $\Omega_D$ can be written uniquely as a sum
$\sum f_i\text{d}x_i$ with for all $e$ only finitely many $f_i$
not in $p^eD$. Moreover, the maps
$\text{d}_{D_e/A, \bar\gamma} : D_e \to \Omega_{D_e/A, \bar\gamma}$
fit together to define a divided power $A$-derivation
\begin{equation}
\label{equation-derivation-D}
\text{d} : D \longrightarrow \Omega_D
\end{equation}
on $p$-adic completions.

\medskip\noindent
We will also need the ``products $\Spec(D(n))$ of $\Spec(D)$'', see
Proposition \href{https://stacks.math.columbia.edu/tag/07JN}{proposition compute cohomology (Tag 07JN)} and its proof for an
explanation. Formally these are defined as follows. For $n \geq 0$ let
$J(n) = \Ker(P \otimes_A \ldots \otimes_A P \to C)$ where
the tensor product has $n + 1$ factors. We set
\begin{equation}
\label{equation-Dn}
D(n) = \lim_e
D_{P \otimes_A \ldots \otimes_A P, \gamma}(J(n))/
p^eD_{P \otimes_A \ldots \otimes_A P, \gamma}(J(n))
\end{equation}
equal to the $p$-adic completion of the divided power envelope.
We denote $\bar J(n)$ its divided power ideal and $\bar \gamma(n)$
its divided powers. We also introduce $D(n)_e = D(n)/p^eD(n)$ as well
as the $p$-adically completed module of differentials
\begin{equation}
\label{equation-omega-Dn}
\Omega_{D(n)} = \lim_e \Omega_{D(n)_e/A, \bar\gamma} =
\lim_e \Omega_{D(n)/A, \bar\gamma}/p^e\Omega_{D(n)/A, \bar\gamma}
\end{equation}
and derivation
\begin{equation}
\label{equation-derivation-Dn}
\text{d} : D(n) \longrightarrow \Omega_{D(n)}
\end{equation}
Of course we have $D = D(0)$. Note that the rings $D(0), D(1), D(2), \ldots$
form a cosimplicial object in the category of divided power rings.



\noindent
In the lemma below we will consider pairs $(M, \nabla)$ satisfying the
following conditions
\begin{enumerate}
\item
\label{item-complete}
$M$ is a $p$-adically complete $D$-module,
\item
\label{item-connection}
$\nabla : M \to M \otimes^\wedge_D \Omega_D$ is a connection, i.e.,
$\nabla(fm) = m \otimes \text{d}f + f\nabla(m)$,
\item
\label{item-integrable}
$\nabla$ is integrable
(see Remark \href{https://stacks.math.columbia.edu/tag/07I0}{remark connection (Tag 07I0)}), and
\item
\label{item-topologically-quasi-nilpotent}
$\nabla$ is {\it topologically quasi-nilpotent}: If we write
$\nabla(m) = \sum \theta_i(m)\text{d}x_i$ for some operators
$\theta_i : M \to M$, then for any $m \in M$ there are only finitely
many pairs $(i, k)$ such that $\theta_i^k(m) \not \in pM$.
\end{enumerate}
The operators $\theta_i$ are sometimes denoted
$\nabla_{\partial/\partial x_i}$ in the literature.
In the following lemma we construct a functor from crystals in quasi-coherent
modules on $\text{Cris}(X/S)$ to the category of such pairs. We will show
this functor is an equivalence in
Proposition \href{https://stacks.math.columbia.edu/tag/07JH}{proposition crystals on affine (Tag 07JH)}.



\begin{situation}
\label{situation-affine}
Here $p$ is a prime number, $(A, I, \gamma)$ is a divided power
ring such that $A$ is a $\mathbf{Z}_{(p)}$-algebra, and $A \to C$ is a
ring map such that $IC = 0$ and such that $p$ is nilpotent in $C$.
\end{situation}

\begin{lemma}
\label{lemma-crystals-on-affine}
In the situation above there is a functor
$$
\begin{matrix}
\text{crystals in quasi-coherent} \\
\mathcal{O}_{X/S}\text{-modules on }\text{Cris}(X/S)
\end{matrix}
\longrightarrow
\begin{matrix}
\text{pairs }(M, \nabla)\text{ satisfying} \\
\text{(\ref{item-complete}), (\ref{item-connection}),
(\ref{item-integrable}), and (\ref{item-topologically-quasi-nilpotent})}
\end{matrix}
$$
\end{lemma}

\begin{proof}
Let $\mathcal{F}$ be a crystal in quasi-coherent modules on $X/S$.
Set $T_e = \Spec(D_e)$ so that $(X, T_e, \bar\gamma)$ is an object
of $\text{Cris}(X/S)$ for $e \gg 0$. We have morphisms
$$
(X, T_e, \bar\gamma) \to (X, T_{e + 1}, \bar\gamma) \to \ldots
$$
which are closed immersions. We set
$$
M =
\lim_e \Gamma((X, T_e, \bar\gamma), \mathcal{F}) =
\lim_e \Gamma(T_e, \mathcal{F}_{T_e}) = \lim_e M_e
$$
Note that since $\mathcal{F}$ is locally quasi-coherent we have
$\mathcal{F}_{T_e} = \widetilde{M_e}$. Since $\mathcal{F}$ is a
crystal we have $M_e = M_{e + 1}/p^eM_{e + 1}$. Hence we see that
$M_e = M/p^eM$ and that $M$ is $p$-adically complete, see
Algebra, Lemma \href{https://stacks.math.columbia.edu/tag/09B8}{lemma limit complete (Tag 09B8)}.

\medskip\noindent
By Lemma \href{https://stacks.math.columbia.edu/tag/07J6}{lemma automatic connection (Tag 07J6)} we know that $\mathcal{F}$
comes endowed with a canonical integrable connection
$\nabla : \mathcal{F} \to \mathcal{F} \otimes \Omega_{X/S}$.
If we evaluate this connection on the objects $T_e$ constructed above
we obtain a canonical integrable connection
$$
\nabla : M \longrightarrow M \otimes^\wedge_D \Omega_D
$$
To see that this is topologically nilpotent we work out what this means.

\medskip\noindent
Now we can do the same procedure for the rings $D(n)$.
This produces a $p$-adically complete $D(n)$-module $M(n)$. Again using
the crystal property of $\mathcal{F}$ we obtain isomorphisms
$$
M \otimes^\wedge_{D, p_0} D(1) \rightarrow M(1)
\leftarrow M \otimes^\wedge_{D, p_1} D(1)
$$
compare with the proof of Lemma \href{https://stacks.math.columbia.edu/tag/07J6}{lemma automatic connection (Tag 07J6)}.
Denote $c$ the composition from left to right. Pick $m \in M$.
Write $\xi_i = x_i \otimes 1 - 1 \otimes x_i$.
Using (\href{https://stacks.math.columbia.edu/tag/07L3}{equation structure Dn (Tag 07L3)}) we can write uniquely
$$
c(m \otimes 1) = \sum\nolimits_K \theta_K(m) \otimes \prod \xi_i^{[k_i]}
$$
for some $\theta_K(m) \in M$ where the sum is over multi-indices
$K = (k_i)$ with $k_i \geq 0$ and $\sum k_i < \infty$. Set
$\theta_i = \theta_K$ where $K$ has a $1$ in the $i$th spot and
zeros elsewhere. We have
$$
\nabla(m) = \sum \theta_i(m) \text{d}x_i.
$$
as can be seen by comparing with the definition of
$\nabla$. Namely, the defining equation is
$p_1^*m = \nabla(m) - c(p_0^*m)$ in Lemma \href{https://stacks.math.columbia.edu/tag/07J6}{lemma automatic connection (Tag 07J6)}
but the sign works out because in the Stacks project we consistently use
$\text{d}f = p_1(f) - p_0(f)$ modulo the ideal of the diagonal squared,
and hence $\xi_i = x_i \otimes 1 - 1 \otimes x_i$ maps to $-\text{d}x_i$
modulo the ideal of the diagonal squared.

\medskip\noindent
Denote $q_i : D \to D(2)$ and $q_{ij} : D(1) \to D(2)$ the coprojections
corresponding to the indices $i, j$. As in the last paragraph of the proof of
Lemma \href{https://stacks.math.columbia.edu/tag/07J6}{lemma automatic connection (Tag 07J6)}
we see that
$$
q_{02}^*c = q_{12}^*c \circ q_{01}^*c.
$$
This means that
$$
\sum\nolimits_{K''} \theta_{K''}(m) \otimes \prod {\zeta''_i}^{[k''_i]}
=
\sum\nolimits_{K', K} \theta_{K'}(\theta_K(m))
\otimes \prod {\zeta'_i}^{[k'_i]} \prod \zeta_i^{[k_i]}
$$
in $M \otimes^\wedge_{D, q_2} D(2)$ where
\begin{align*}
\zeta_i & = x_i \otimes 1 \otimes 1 - 1 \otimes x_i \otimes 1,\\
\zeta'_i & = 1 \otimes x_i \otimes 1 - 1 \otimes 1 \otimes x_i,\\
\zeta''_i & = x_i \otimes 1 \otimes 1 - 1 \otimes 1 \otimes x_i.
\end{align*}
In particular $\zeta''_i = \zeta_i + \zeta'_i$ and we have that
$D(2)$ is the $p$-adic completion of the divided power polynomial
ring in $\zeta_i, \zeta'_i$ over $q_2(D)$, see Lemma \href{https://stacks.math.columbia.edu/tag/07L2}{lemma structure Dn (Tag 07L2)}.
Comparing coefficients in the expression above it follows immediately that
$\theta_i \circ \theta_j = \theta_j \circ \theta_i$
(this provides an alternative proof of the integrability of $\nabla$) and that
$$
\theta_K(m) = (\prod \theta_i^{k_i})(m).
$$
In particular, as the sum expressing $c(m \otimes 1)$ above has to converge
$p$-adically we conclude that for each $i$ and each $m \in M$ only a finite
number of $\theta_i^k(m)$ are allowed to be nonzero modulo $p$.
\end{proof}
\end{document}
